\documentclass[11pt]{article}
\usepackage[margin=1in]{geometry}
\usepackage{amsmath,amssymb,amsthm}
\usepackage[T1]{fontenc}
\usepackage{lmodern}
\usepackage[hidelinks]{hyperref}
\newtheorem{theorem}{Theorem}
\newtheorem{corollary}[theorem]{Corollary}
\title{Conjecture 00000001259:\\The Tribonacci matrix does not have three real eigenvalues}
\author{C0ldSmi1e}
\date{4 October 2026}
\begin{document}
\maketitle
\begin{abstract}
The substitution $a\mapsto ab$, $b\mapsto ac$, $c\mapsto a$ has incidence matrix with characteristic polynomial $X^3-X^2-X-1$ and discriminant $-44$. Its characteristic polynomial therefore does not split over the real numbers, contradicting the conjecture's explicit assertion of three real eigenvalues, including when eigenvalues are counted with algebraic multiplicity. The Lean proof constructs the matrix from letter counts and uses its actual characteristic polynomial.
\end{abstract}

\section{The assertion being disproved}
The repository's English statement is:
\begin{quote}
Definition: The Tribonacci word is the fixed point of the substitution $a\to ab$, $b\to ac$, $c\to a$; the spectral hull of its substitution matrix is the algebraic hull of the spectrum. Conjecture: The spectral hull has dimension 3 (generated by three real eigenvalues); the second eigenvector corresponds to the second-largest eigenvalue and gives an algebraic characterization of the combinatorial correction coefficient $0.54\ldots$ of the word-frequency vector.
\end{quote}
The exact bilingual source is preserved in \texttt{conjecture.md}. Both language versions explicitly assert that the matrix has three real eigenvalues. We refute this necessary part of the statement. We do not supply a new definition of ``spectral hull'' or ``algebraic hull,'' and do not separately formalize its dimension or the correction coefficient.

Let the ordered alphabet be $(a,b,c)$ and let $\sigma$ denote the given substitution. The incidence convention is
\[
 M_{ij}=\text{the number of occurrences of letter }i\text{ in }\sigma(j).
\]
Counting letters in the three substituted words gives
\[
 M=\begin{pmatrix}1&1&1\\1&0&0\\0&1&0\end{pmatrix}.
\]
The alternative convention counting in rows gives $M^{\mathsf T}$, which has the same characteristic polynomial. No construction of the infinite fixed word is required: the conjectured spectral assertion concerns this substitution matrix, determined by the three finite words.

\section{A discriminant obstruction}
\begin{theorem}
The characteristic polynomial of $M$ does not split over $\mathbb R$.
\end{theorem}
\begin{proof}
Direct determinant expansion gives
\[
 \chi_M(X)=\det(XI-M)
 =\det\begin{pmatrix}X-1&-1&-1\\-1&X&0\\0&-1&X\end{pmatrix}
 =X^3-X^2-X-1.
\]
For a cubic $aX^3+bX^2+cX+d$, its discriminant is
\[
 \Delta=b^2c^2-4ac^3-4b^3d-27a^2d^2+18abcd.
\]
Substituting $(a,b,c,d)=(1,-1,-1,-1)$ yields
\[
 \Delta=1+4-4-27-18=-44.
\]
If $\chi_M$ split over $\mathbb R$, its three roots, counted with multiplicity, could be written $x,y,z\in\mathbb R$. The standard discriminant identity for a monic cubic then gives
\[
 -44=\Delta=((x-y)(x-z)(y-z))^2\geq0,
\]
a contradiction. This argument allows repeated roots; if two roots coincide the square is zero, which is still incompatible with $-44$.
\end{proof}
\begin{corollary}
The explicit assertion of three real eigenvalues in conjecture 00000001259 is false, for either incidence-matrix convention.
\end{corollary}
\begin{proof}
Eigenvalues counted with algebraic multiplicity are the roots of the characteristic polynomial. A real $3\times3$ matrix with three real eigenvalues, counted with that multiplicity, has a characteristic polynomial that splits over $\mathbb R$. The theorem rules this out for $M$, and $M^{\mathsf T}$ has the same characteristic polynomial. In particular, three distinct real eigenvalues are also impossible.
\end{proof}

\clearpage
\section{Formalization and verification}
The Lean development uses \texttt{Fin 3} for the ordered alphabet, with $0=a$, $1=b$, and $2=c$. The substitution is defined by the actual lists $[0,1]$, $[0,2]$, and $[0]$. Each incidence entry is a \texttt{List.count}, cast to $\mathbb R$. Equality with the displayed matrix is proved before the characteristic polynomial is computed with Mathlib's \texttt{Matrix.charpoly}.

The cubic is Mathlib's \texttt{Cubic} with coefficients $(1,-1,-1,-1)$. The proof verifies its polynomial and discriminant exactly. Mathlib's splitting theorem supplies a multiset of three real roots, permitting repetition; its discriminant identity gives the real square above. Thus no unproved numerical root approximation or external classification of cubic roots enters the argument.

The formal predicate \texttt{threeRealEigenvalues} asserts that the multiset of real roots of the actual characteristic polynomial has cardinality three. The theorem \texttt{threeRealEigenvalues\_iff\_splits} proves that this is equivalent to splitting over $\mathbb R$, using the matrix dimension to establish degree three. The final theorem \texttt{conjecture1259\_disproof} negates that predicate. It is not a substitute definition of the undefined hull terminology. No separate claim is made here about the hull's dimension, a second eigenvector, the coefficient $0.54\ldots$, or a formal count of distinct real roots.

The project pins Lean 4.19.0 and Mathlib v4.19.0, with the exact dependency revisions in \texttt{lean/lake-manifest.json}. From \texttt{lean/}, run:
\begin{verbatim}
lake build
lake env lean -DwarningAsError=true Conjecture1259.lean
lake env lean -DwarningAsError=true Check.lean
\end{verbatim}
The audit file prints the definitions, theorem types, and axiom dependencies. The submission includes local verification records and independent semantic scrutiny. No auxiliary numerical program is used. Local validation does not constitute maintainer acceptance.

\paragraph{Eligibility and source.}
At the recorded pre-submission check, the repository metadata marked this conjecture unsolved and no earlier submission for it was identified in the solution history or pull requests. The exact check details are in \texttt{verification/eligibility.json}. The conjecture text and contribution requirements were read at repository revision \texttt{0862407e}. The source is available at \href{https://github.com/The-Last-Math-Competition/The-Last-Math-Competition/blob/0862407ef50dda4f7376342ca3e79368dce942d2/conjectures/00000001259.md}{the immutable conjecture page}.
\end{document}
